% =============================================================================
% 04-system-programs-os-structure-notes.tex — Lecture Notes for Lecture 4:
% System Programs and Operating-System Structures
%
% DERIVED ARTIFACT. Source of truth: Slides/CS502/04-system-programs-os-structure.tex.
% Generated by /lecture-notes; do not hand-edit content independently of the
% Beamer source without re-running /qa-notes afterward (single-source-of-truth.md).
% Structured as a textbook chapter (numbered sections/figures/examples,
% topic-organized rather than a frame-by-frame transcript). TikZ figures are
% copied verbatim from the deck.
%
% Anchor reading: Silberschatz, Galvin & Gagne, Operating System Concepts,
%                 8th ed. — Ch. 2 §2.5 (p. 66), §2.6, §2.7 (pp. 70-76), §2.8
%                 (p. 76); Stallings 6e Ch. 2 §2.3 (p. 64), Ch. 4 §4.3 (p. 179);
%                 Bovet & Cesati 3e Ch. 1 (pp. 1-34). Page-verified via
%                 master_supporting_docs/CS502/supporting_books/*/index.md.
%
% Compile with: cd Notes/CS502 &&
%   TEXINPUTS=../../Preambles:$TEXINPUTS xelatex -interaction=nonstopmode 04-system-programs-os-structure-notes.tex  (x3)
% =============================================================================

\documentclass[11pt]{article}
\usepackage[margin=1in]{geometry}
\usepackage{array}
\input{../../Preambles/header}
\coursecode{CS502}

\renewcommand{\thesection}{4.\arabic{section}}
\renewcommand{\thefigure}{4.\arabic{figure}}
\renewcommand{\thetable}{4.\arabic{table}}

\newtheorem{example}{Example}
\renewcommand{\theexample}{4.\arabic{example}}

\newenvironment{definitionbox}[1]{%
  \par\noindent\textbf{Definition (#1).}\ \itshape
}{\par\vspace{0.3em}}

\setlength{\emergencystretch}{2em}

\title{System Programs and Operating-System Structures\\\large Lecture Notes --- Lecture 4\\
  \normalsize Monolithic, Layered, Microkernel, Modular, Virtual Machine}
\author{Auqib Hamid Lone\\[0.3em]
  \emph{Department of Computer Science \& Engineering}, \emph{GCET Ganderbal}}
\date{\today}

\begin{document}
\maketitle

\section{One Bad Driver}

Imagine a network-card driver with a bug: on a rare packet, it dereferences a
\texttt{NULL} pointer. Does the whole machine crash, does only the network stop,
or does nobody notice? The answer is \emph{not} in the driver's code. It is in
how the operating system is \key{organized}. Where the driver runs --- in kernel
mode or user mode, inside the kernel or beside it --- decides how far the damage
spreads. This chapter drops that same buggy driver into five structures in turn
and watches its \emph{blast radius} change.

Lecture~3 studied the system-call doorway. This chapter studies what sits on each
side of it. Above the doorway are the \key{system programs}, ordinary user-mode
programs that make system calls. Below it is the kernel, whose internal
organization --- simple, monolithic, layered, microkernel, modular, or a guest
of a virtual machine --- is its \key{structure}. Between the two we need a
small amount of design vocabulary: design goals, and the separation of
mechanism from policy. The chapter closes with a single comparison table: one
driver bug, five outcomes.

\noindent\muted{\textit{Reading: Silberschatz et al., Ch.~2 \S2.5--2.8,
pp.~66--84 \cite{SilberschatzOSC_operating_system_concepts}; Stallings, Ch.~2
and Ch.~4 \S4.3 \cite{StallingsOS_internals_design_principles}.}}

\section{System Programs}

\noindent\muted{\textit{Reading: Silberschatz et al., Ch.~2 \S2.5, p.~66.}}

Most of what a user calls ``the operating system'' is not the kernel at all. It
is a set of ordinary programs that make system calls.

\begin{definitionbox}{System program}
A \key{system program} (system utility) is a program shipped with the operating
system that gives users a convenient environment for program development and
execution. It runs in user mode and obtains every privileged service through
system calls.
\end{definitionbox}

\texttt{ls}, \texttt{cp}, \texttt{cc}, \texttt{vi}, the shell, and the login
program are all system programs. Some are thin wrappers around one system
call; others are large programs in their own right. Either way, the user's view
of the operating system is defined by the system programs, not by the system
calls: a user types \texttt{ls}, not \texttt{getdents}.

Silberschatz groups system programs into six categories
(Table~\ref{tab:sysprog}). Beyond these, many systems also ship
\key{application programs} --- browsers, word processors, databases --- which
are not part of the operating system proper.

\begin{table}[h]
\centering
\small
\renewcommand{\arraystretch}{1.3}
\begin{tabular}{@{}>{\raggedright\arraybackslash}p{0.27\linewidth}>{\raggedright\arraybackslash}p{0.40\linewidth}>{\raggedright\arraybackslash}p{0.25\linewidth}@{}}
\toprule
\textbf{Category} & \textbf{What it does} & \textbf{Examples} \\
\midrule
File management & create, delete, copy, rename, list files and directories & \texttt{cp}, \texttt{mv}, \texttt{rm}, \texttt{ls} \\
Status information & report date, time, memory, users, performance & \texttt{date}, \texttt{who}, \texttt{df} \\
File modification & create and change file contents & \texttt{vi}, \texttt{grep}, \texttt{sed} \\
Programming-language support & compilers, assemblers, interpreters, debuggers & \texttt{cc}, \texttt{as}, \texttt{gdb} \\
Program loading and execution & loaders, linkage editors & \texttt{ld}, the loader \\
Communications & connections between processes, users, and machines & \texttt{ssh}, \texttt{mail}, \texttt{ftp} \\
\bottomrule
\end{tabular}
\caption{The categories of system programs (Silberschatz et al., Ch.~2 \S2.5,
p.~66).}
\label{tab:sysprog}
\end{table}

\begin{example}[Classifying six commands]
Classify each command by system-program category.
\begin{enumerate}
  \item \texttt{ls -l} lists a directory's entries, using \texttt{open},
        \texttt{getdents}-style reads, and \texttt{stat}: \textbf{file
        management}.
  \item \texttt{date} reads the clock through an information-maintenance
        system call: \textbf{status information}.
  \item \texttt{vi notes.txt} edits contents with \texttt{read} and
        \texttt{write}: \textbf{file modification}.
  \item \texttt{cc prog.c} compiles source into an object file:
        \textbf{programming-language support}.
  \item \texttt{ld prog.o} links objects into an executable: \textbf{program
        loading and execution}.
  \item \texttt{ssh server} opens a connection with \texttt{socket}:
        \textbf{communications}.
\end{enumerate}
Notice that there are two taxonomies in play. A system \emph{program}
category (this chapter) is not the same thing as a system \emph{call} category
(Lecture~3). \texttt{date} is a status-information \emph{program} that makes an
information-maintenance \emph{call}.
\end{example}

\subsection{The command interpreter}

The \key{shell} reads a command and gets it executed. There are two ways to
build it (Silberschatz et al., Ch.~2 \S2.2). In the first, the interpreter
contains the code for each command as a \emph{built-in}; adding a command means
changing the interpreter, and the interpreter grows with every command. In the
second, used by UNIX, commands are \emph{system programs}: \texttt{rm file.txt}
makes the shell find a \emph{file} named \texttt{rm}, load it, and pass
\texttt{file.txt} as its argument. New commands are new files, and the shell
stays small. The UNIX shell uses the second design for most commands, and keeps
a few built-ins such as \texttt{cd}.

Why must \texttt{cd} be a built-in? A separate program would run in a
\emph{child} process created by \texttt{fork}. Calling \texttt{chdir} in that
child changes only the child's current directory, and the child then exits.
The shell's own directory would be unchanged, so the shell must call
\texttt{chdir} itself.

\section{Design Principles}

\noindent\muted{\textit{Reading: Silberschatz et al., Ch.~2 \S2.6.}}

Before choosing a structure, a designer decides what the system is for. Users
want a system that is convenient, easy to learn, reliable, safe, and fast
(\key{user goals}). The people who build it want one that is easy to design,
implement, and maintain, and that is flexible, reliable, error-free, and
efficient (\key{system goals}). These goals conflict, so there is no single
``best'' operating system --- only a best fit to a set of requirements.

The second principle keeps the decisions about \emph{what} to do apart from the
machinery for \emph{how} to do it.

\begin{definitionbox}{Mechanism vs.\ policy}
A \key{mechanism} determines \emph{how} to do something; a \key{policy}
determines \emph{what} will be done. Separating them lets a policy change
without rewriting the mechanism.
\end{definitionbox}

\begin{example}[Separating mechanism from policy]
Table~\ref{tab:mechpol} separates four situations into their mechanism and
policy.
\begin{table}[h]
\centering
\small
\renewcommand{\arraystretch}{1.3}
\begin{tabular}{@{}>{\raggedright\arraybackslash}p{0.20\linewidth}>{\raggedright\arraybackslash}p{0.36\linewidth}>{\raggedright\arraybackslash}p{0.36\linewidth}@{}}
\toprule
\textbf{Situation} & \textbf{Mechanism (how)} & \textbf{Policy (what)} \\
\midrule
CPU protection & the timer interrupts the CPU after a set interval & how long the interval is for each user \\
Scheduling & a context switch from one process to another & which process runs next \\
File protection & the kernel checks a permission bit before \texttt{open} & who is granted which bits \\
Our driver & loading code into the kernel at run time & which drivers are trusted enough to load \\
\bottomrule
\end{tabular}
\caption{Mechanism versus policy in four situations.}
\label{tab:mechpol}
\end{table}
The split pays off because the timer mechanism of Lecture~2 is unchanged whether
the quantum is 10\,ms or 100\,ms; only a parameter changes. A microkernel
(Section~\ref{sec:micro}) pushes this idea to its limit: the kernel keeps the
mechanisms, and user-level servers hold the policy.
\end{example}

\subsection{Implementation language}

Early operating systems were written in \key{assembly language}. Modern kernels
such as Linux and Windows are written mostly in \key{C}, with small parts in
assembly: the trap entry, the context switch, and device-register access. A
high-level language is faster to write, easier to read and debug, and easier to
\key{port} to new hardware. The possible cost is slower code, but the
bottlenecks are usually algorithms and data structures, which can be tuned
after profiling. The structure of the code therefore matters more than its
language: a well-structured C kernel can be optimized where it counts.

\section{Simple and Monolithic Structure}

\noindent\muted{\textit{Reading: Silberschatz et al., Ch.~2 \S2.7,
pp.~70--76.}}

A kernel is millions of lines. How those lines are divided into parts --- and
which parts share kernel mode --- is its structure.

\subsection{Simple structure: MS-DOS}

MS-DOS was written for small machines, and its levels are not well separated
(Figure~\ref{fig:msdos}). Applications can call BIOS routines and write to the
display directly, because the Intel 8088 for which it was written had \key{no
dual mode} and no hardware protection. The consequence is that one bad program
can crash the whole system.

\begin{figure}[h]
\centering
\begin{tikzpicture}[every node/.style={font=\small},
      lvl/.style={draw=primary-blue, fill=light-bg, minimum width=6.4cm,
                  minimum height=0.62cm, align=center}]
    \node[lvl] (app)  at (0, 1.65) {Application program};
    \node[lvl] (res)  at (0, 0.85) {Resident system program};
    \node[lvl] (dos)  at (0, 0.05) {MS-DOS device drivers};
    \node[lvl, draw=primary-gold, fill=white] (bios) at (0,-0.75)
      {ROM BIOS device drivers};
    \draw[confound-edge] (app.east) -- (4.0,1.65) |- (bios.east);
    \node[anchor=west, text=negative, align=left, font=\footnotesize]
      at (4.25,0.45) {direct access\\to hardware};
  \end{tikzpicture}
\caption{MS-DOS layer structure. The dashed path shows an application bypassing
the middle levels to reach the hardware directly --- cf.\ Silberschatz et al.,
Ch.~2 \S2.7, pp.~70--76 \cite{SilberschatzOSC_operating_system_concepts}.}
\label{fig:msdos}
\end{figure}

\subsection{Monolithic structure: traditional UNIX}

Traditional UNIX has two parts: the \key{kernel} and the \key{system programs}
(Figure~\ref{fig:unix}). Everything between the system-call interface and the
hardware interface --- file system, CPU scheduling, memory management, signals,
terminal handling, and device drivers --- is \emph{one} program running in
kernel mode in one address space. That makes it fast, since a request from one
subsystem to another is just a function call, but hard to maintain.

\begin{figure}[h]
\centering
\begin{tikzpicture}[every node/.style={font=\footnotesize},
      bx/.style={draw=primary-blue, fill=light-bg, minimum width=9.6cm,
                 align=center},
      bd/.style={draw=negative, dashed, fill=white, minimum width=9.6cm,
                 align=center}]
    \node[bx, minimum height=0.50cm] (usr) at (0, 1.75)
      {Users \quad shells and commands};
    \node[bx, minimum height=0.50cm] (sys) at (0, 1.15)
      {System programs: compilers, interpreters, system libraries};
    \node[bd, minimum height=0.44cm] (sci) at (0, 0.58)
      {System-call interface to the kernel};
    \node[bx, minimum height=1.00cm, fill=white, draw=primary-gold, text width=9.2cm] (krn) at (0,-0.25)
      {\textbf{Kernel:} file system, CPU scheduling, memory management,
      signals, terminal handling, \key{device drivers} --- one address space};
    \node[bd, minimum height=0.44cm] (hwi) at (0,-1.07)
      {Kernel interface to the hardware};
    \node[bx, minimum height=0.50cm] (hw) at (0,-1.64)
      {Hardware: terminal controllers, device controllers, memory};
  \end{tikzpicture}
\caption{Traditional UNIX system structure: one large kernel between two
interfaces --- cf.\ Silberschatz et al., Ch.~2 \S2.7, pp.~70--76
\cite{SilberschatzOSC_operating_system_concepts}.}
\label{fig:unix}
\end{figure}

\begin{example}[The bad driver in a monolithic kernel]
Follow the buggy driver step by step.
\begin{enumerate}
  \item The driver is compiled into the kernel and runs in \key{kernel mode}
        in the kernel's address space.
  \item The rare packet arrives, and the driver dereferences \texttt{NULL}.
  \item The fault happens in kernel mode, so there is no more-privileged code
        to contain it.
  \item The kernel cannot tell which of its data structures the driver
        corrupted, so it halts (a \emph{kernel panic}).
\end{enumerate}
\textbf{Outcome:} the blast radius is \bad{the whole machine}; every process,
not just the network, is lost. The same speed that makes monolithic kernels
attractive --- no boundary between subsystems --- is what lets one bug reach
everything.
\end{example}

\section{The Layered Approach}

\noindent\muted{\textit{Reading: Silberschatz et al., Ch.~2 \S2.7,
pp.~70--76; Stallings, Ch.~2 \S2.3, p.~64.}}

In the layered approach the operating system is a stack of \key{layers}, from
layer~0 (the hardware) to layer~$N$ (the user interface), and layer $i$ uses
only the operations of layers below $i$ (Figure~\ref{fig:layers}). Each layer
hides its data structures from the layers above. The great advantage is that
the system can be \good{debugged bottom-up}: once layer $i-1$ has been
verified, any bug found while testing layer $i$ must be in layer $i$.

\begin{figure}[h]
\centering
\begin{tikzpicture}[every node/.style={font=\footnotesize},
          ly/.style={draw=primary-blue, fill=light-bg, minimum width=4.6cm,
                     minimum height=0.50cm, align=center}]
        \node[ly] (L5) at (0, 1.80) {Layer $N$: user interface};
        \node[ly] (L4) at (0, 1.05) {I/O management};
        \node[ly] (L3) at (0, 0.30) {Device drivers, console};
        \node[ly] (L2) at (0,-0.45) {Memory management};
        \node[ly] (L1) at (0,-1.20) {CPU scheduling};
        \node[ly, draw=primary-gold, fill=white] (L0) at (0,-1.95)
          {Layer 0: hardware};
        \foreach \a/\b in {L5/L4, L4/L3, L3/L2, L2/L1, L1/L0}
          \draw[observed-edge] (\a.south) -- (\b.north);
      \end{tikzpicture}
\caption{A layered operating system; each arrow is a call to the layer
immediately below --- cf.\ Silberschatz et al., Ch.~2 \S2.7, pp.~70--76
\cite{SilberschatzOSC_operating_system_concepts}; Stallings, Ch.~2 \S2.3,
p.~64 \cite{StallingsOS_internals_design_principles}.}
\label{fig:layers}
\end{figure}

\begin{example}[Where does the backing-store driver go?]
The rule is: if A calls B, then B must be in a \emph{lower} layer than A.
\begin{enumerate}
  \item Memory management swaps pages to disk, so it calls the backing-store
        driver. Therefore the \textbf{driver is below memory management}.
  \item The driver waits for I/O, and the CPU should run someone else in the
        meantime, so the driver calls the CPU scheduler. Therefore the
        \textbf{scheduler is below the driver}.
  \item But on a large system the scheduler's tables may not fit in memory
        and must themselves be swapped, so the scheduler calls the driver.
        Therefore the \textbf{driver is below the scheduler}.
\end{enumerate}
Steps 2 and 3 demand opposite orders. Choosing the layers requires careful
planning, and some dependencies have no clean order at all.
\end{example}

The example illustrates the first cost of layering: \key{defining the layers is
hard}, and the cycle above is typical. The second cost is that \key{each layer
adds overhead}: a request from the top passes through every layer below it,
each adding a call and possibly copying parameters. A system call that needs
I/O traverses I/O $\to$ memory $\to$ scheduling $\to$ hardware, even when a
direct path would do. Modern systems therefore use fewer, more functional
layers, keeping the benefits without the full cost.

\begin{example}[The bad driver, layered]
Layering organizes code; it does not change privilege. The driver layer still
runs in kernel mode, so the \texttt{NULL} dereference still \bad{crashes the
whole kernel}. The win from layering is finding the bug sooner, not surviving
it.
\end{example}

\section{Microkernels}
\label{sec:micro}

\noindent\muted{\textit{Reading: Silberschatz et al., Ch.~2 \S2.7; Stallings,
Ch.~4 \S4.3, p.~179.}}

If kernel mode is where one bug becomes everyone's bug, the natural response is
to put as little as possible in kernel mode. A \key{microkernel} keeps only
message passing (IPC), basic memory management, and CPU scheduling in the
kernel, and moves everything nonessential --- file systems, drivers, networking
--- into \key{user-level servers} (Figure~\ref{fig:micro}). Clients and servers
never call each other directly: they exchange \key{messages} through the
microkernel. Mach (developed at CMU in the mid-1980s) and QNX are the standard
examples.

\begin{figure}[h]
\centering
\begin{tikzpicture}[every node/.style={font=\footnotesize},
      us/.style={draw=primary-blue, fill=light-bg, minimum width=2.6cm,
                 minimum height=0.70cm, align=center},
      kb/.style={draw=primary-gold, fill=white, minimum width=10.0cm,
                 align=center}]
    \node[us] (app) at (-3.6, 1.30) {Application};
    \node[us] (fs)  at ( 0.0, 1.30) {File server};
    \node[us] (drv) at ( 3.6, 1.30) {Device driver};
    \draw[dashed, draw=negative] (-5.0,0.55) -- (5.0,0.55);
    \node[anchor=west, text=negative] at (5.1,0.75) {user mode};
    \node[anchor=west, text=negative] at (5.1,0.35) {kernel mode};
    \node[kb, minimum height=0.80cm] (mk) at (0,-0.25)
      {\textbf{Microkernel:} message passing (IPC), basic memory management,
      CPU scheduling};
    \node[kb, minimum height=0.50cm, draw=primary-blue, fill=light-bg] (hw)
      at (0,-1.20) {Hardware};
    \draw[observed-edge] (-3.6,0.95) -- (-3.6,0.15);
    \draw[observed-edge] (-0.6,0.15) -- (-0.6,0.95);
    \node[anchor=west, text=primary-gold] at (-3.4,0.75) {send};
    \node[anchor=west, text=primary-gold] at (-0.4,0.75) {deliver};
  \end{tikzpicture}
\caption{Microkernel structure. One message: the application sends to the
microkernel, which delivers to the file server --- cf.\ Silberschatz et al.,
Ch.~2 \S2.7; Stallings, Ch.~4 \S4.3, p.~179
\cite{StallingsOS_internals_design_principles}.}
\label{fig:micro}
\end{figure}

\begin{example}[Counting boundary crossings]
The application reads a file block. Count the traps, where each trap is one
user$\to$kernel entry plus one kernel$\to$user return, i.e.\ 2 crossings.
\begin{enumerate}
  \item \textbf{Monolithic.} One \texttt{read} trap; the file system and driver
        are kernel functions. Result: \textbf{1 trap, 2 crossings}.
  \item \textbf{Microkernel, file server only.} The application sends to the
        server (1 trap) and the server replies to the application (1 trap).
        Result: \textbf{2 traps, 4 crossings}.
  \item \textbf{Microkernel, file server and driver server.} The messages are
        app$\to$fs, fs$\to$drv, drv$\to$fs, fs$\to$app. Result:
        \textbf{4 traps, 8 crossings}.
\end{enumerate}
\textbf{General pattern (our counting model).} A request that passes through a
chain of $k$ servers and back sends $k$ messages on the way out and $k$ replies
on the way back, so it costs $2k$ message traps, i.e.\ $4k$ crossings, versus 2
crossings for a monolithic call. With $k=2$ that is $8/2 = 4$ times as many.
Each message may also be copied between address spaces; this is the
``system-function overhead'' that makes microkernels slower.
\end{example}

The benefits of the microkernel are that it is \good{extensible} (a new service
is a new user-level server and the kernel is untouched), \good{portable}
(little hardware-specific code lives in the kernel), and \good{reliable and
secure} (most services run as user processes). Stallings, Ch.~4 \S4.3,
p.~179, lists further benefits: uniform interfaces, flexibility, and support
for distributed systems. The cost is \bad{performance}: messages and extra
crossings for every service request.

\begin{example}[The bad driver in a microkernel]
The driver is a \emph{user process}. The \texttt{NULL} dereference causes an
ordinary page fault; the kernel terminates that process only, and can restart
it. The blast radius is \good{the network driver}; other processes keep
running.
\end{example}

The history of Windows NT shows the trade-off in practice (Silberschatz et al.,
Ch.~2 \S2.7). NT began with a layered microkernel design and later moved more
services back into the kernel. Its designers were solving the cost of the
microkernel design: every service request had become messages and extra mode
switches. Moving services into kernel space turns messages back into function
calls, which is faster --- but the price is the isolation just described: a bug
in a moved service can again crash the kernel.

\section{Modules and Hybrid Systems}

\noindent\muted{\textit{Reading: Silberschatz et al., Ch.~2 \S2.7, pp.~70--76;
Bovet \& Cesati, Ch.~1, pp.~1--34.}}

A modular kernel consists of a small \key{core kernel} plus \key{loadable
modules} that are linked in at boot time or at run time, without recompiling
the kernel. Each module has a defined interface, like a layer; any module may
call any other, like a microkernel's servers; yet no messages are needed,
because all modules run in kernel mode. Solaris is organized as a core kernel
with seven types of loadable modules: scheduling classes, file systems,
loadable system calls, executable formats, STREAMS modules, miscellaneous
modules, and device and bus drivers. Linux is monolithic with loadable modules
\cite{BovetCesati_understanding_linux_kernel}.

Real kernels mix structures to balance speed against isolation. \key{Linux} is
monolithic plus loadable modules. \key{Mac OS X} combines a Mach microkernel
with a BSD kernel environment, plus loadable kernel extensions.
\key{Windows} is largely monolithic for performance, with microkernel-style
user-mode subsystems.

\begin{example}[The bad driver as a module]
\texttt{insmod netdrv.ko} links the driver into the running kernel, in
\key{kernel mode}. The \texttt{NULL} dereference therefore \bad{crashes the
kernel}, exactly as in the monolithic case. Modules buy \emph{flexibility}
(load, unload, and update without a reboot), not \emph{isolation}.
\end{example}

\section{Virtual Machines}

\noindent\muted{\textit{Reading: Silberschatz et al., Ch.~2 \S2.8, p.~76.}}

Pushing the isolation idea one level further, we can run the entire operating
system as a guest, on hardware that is itself simulated. The
\key{virtual-machine manager} (hypervisor) gives each guest the illusion of its
own processor, memory, and devices (Figure~\ref{fig:vm}), and each guest runs a
full, unmodified operating system --- possibly a different one in each virtual
machine. IBM's VM/370 (1972), VMware, and the Java Virtual Machine are the
standard examples.

\begin{figure}[h]
\centering
\begin{tikzpicture}[every node/.style={font=\footnotesize},
      gb/.style={draw=primary-blue, fill=light-bg, minimum width=3.0cm,
                 minimum height=0.55cm, align=center},
      wd/.style={draw=primary-gold, fill=white, minimum width=9.6cm,
                 align=center}]
    \node[gb] (p1) at (-3.25, 0.60) {processes};
    \node[gb] (p2) at ( 0.00, 0.60) {processes};
    \node[gb] (p3) at ( 3.25, 0.60) {processes};
    \node[gb, fill=white] (k1) at (-3.25,-0.10) {kernel (VM 1)};
    \node[gb, fill=white] (k2) at ( 0.00,-0.10) {kernel (VM 2)};
    \node[gb, fill=white] (k3) at ( 3.25,-0.10) {kernel (VM 3)};
    \node[wd, minimum height=0.55cm] (vmm) at (0,-0.85)
      {Virtual-machine manager (hypervisor)};
    \node[wd, minimum height=0.50cm, draw=primary-blue, fill=light-bg] (hw)
      at (0,-1.55) {Hardware};
  \end{tikzpicture}
\caption{Three virtual machines on one virtual-machine manager --- cf.\
Silberschatz et al., Ch.~2 \S2.8, p.~76
\cite{SilberschatzOSC_operating_system_concepts}.}
\label{fig:vm}
\end{figure}

Virtual machines bring \good{isolation} (each VM is protected from the others,
so a virus in one cannot reach another directly), support for \good{system
development} (a new kernel can be tested in a VM without disrupting normal work
on the same machine), and \good{consolidation} (several lightly used servers
become VMs on one physical machine). The \bad{cost} is that the guest kernel
itself runs in a de-privileged mode, so its privileged operations must be
trapped and emulated by the manager.

\begin{example}[The bad driver in a virtual machine]
The driver runs in the \emph{guest} kernel of VM~1. The guest panics, so
\bad{all of VM~1 is lost}, but the manager and VMs 2 and 3 \good{keep running}.
The blast radius is one virtual machine.
\end{example}

\section{Comparison}

Table~\ref{tab:compare} gathers the five outcomes of the running example. The
pattern is that isolation is bought with boundary crossings, and every
structure picks a point on that trade-off.

\begin{table}[h]
\centering
\small
\renewcommand{\arraystretch}{1.3}
\begin{tabular}{@{}>{\raggedright\arraybackslash}p{0.14\linewidth}>{\raggedright\arraybackslash}p{0.15\linewidth}>{\raggedright\arraybackslash}p{0.17\linewidth}>{\raggedright\arraybackslash}p{0.18\linewidth}>{\raggedright\arraybackslash}p{0.20\linewidth}@{}}
\toprule
\textbf{Structure} & \textbf{Driver runs in} & \textbf{Bug's blast radius} & \textbf{Speed} & \textbf{Main strength} \\
\midrule
Monolithic & kernel mode & whole machine & fastest (calls) & performance \\
Layered & kernel mode & whole machine & layer overhead & bottom-up debugging \\
Microkernel & user mode & the driver only & messages, $4k$ crossings & isolation, extensibility \\
Modular & kernel mode & whole machine & near monolithic & load at run time \\
Virtual machine & guest kernel & one VM & trap-and-emulate & isolates whole OSes \\
\bottomrule
\end{tabular}
\caption{One driver bug, five structures.}
\label{tab:compare}
\end{table}

\begin{example}[Choosing a structure for a flight-control computer]
An aircraft flight-control computer must keep running even if its
network-logging driver fails. A \key{microkernel} is the natural choice: the
logging driver runs as a user-level server, so its failure cannot take down the
control software. The price is that every driver request is a message --- with
a chain of $k$ servers, $4k$ crossings instead of 2. This is why microkernels
such as QNX are common in embedded and real-time systems.
\end{example}

\subsection*{Common misconceptions}

\begin{itemize}
  \item ``The shell is part of the kernel.'' --- No: the shell is a system
        program running in user mode.
  \item ``Layered means the upper layers run in user mode.'' --- No: layering
        organizes code; privilege is unchanged.
  \item ``Loadable modules isolate drivers.'' --- No: modules run in kernel
        mode, and a bug crashes the kernel.
  \item ``Microkernels are faster because the kernel is small.'' --- No:
        messages add crossings; microkernels trade speed for isolation.
\end{itemize}

\subsection*{Summary}

System programs (file management, status information, file modification,
language support, program loading and execution, communications) run in user
mode on top of system calls. Designers separate \key{mechanism} (how) from
\key{policy} (what). A \key{monolithic} kernel is fast and lives in one address
space, so one bug crashes all; a \key{layered} kernel is easy to debug but hard
to order and adds overhead; a \key{microkernel} keeps a minimal kernel and
runs services as user-level servers, paying for isolation in message
crossings; a \key{modular} kernel gains run-time flexibility but is still kernel
mode; and a \key{virtual machine} isolates entire operating systems.

The next lecture turns to the unit every one of these structures manages --- the
\key{process}, and the kernel record that describes it, the \key{process
control block}. The reading is Silberschatz et al., Ch.~3 \S3.1
\cite{SilberschatzOSC_operating_system_concepts}, and Stallings, Ch.~3
\S3.1--3.3 \cite{StallingsOS_internals_design_principles}.

For this chapter the page-cited sources are: Silberschatz et al., 8th ed.,
Ch.~2 \S2.5 (system programs, p.~66), \S2.6 (design and implementation),
\S2.7 (OS structure, pp.~70--76), and \S2.8 (virtual machines, p.~76), with
\S2.5--2.8 spanning pp.~66--84; Stallings, 6th ed., Ch.~2 \S2.3 (system
structure, p.~64) and Ch.~4 \S4.3 (microkernels, p.~179); and Bovet \& Cesati,
3rd ed., Ch.~1 (kernel architecture: monolithic with loadable modules,
pp.~1--34). All are page-verified through each book's \texttt{index.md} under
\texttt{master\_supporting\_docs/CS502/supporting\_books/}.

\section*{Self-Check}

\subsection*{Exercises}

\begin{enumerate}
  \item Classify \texttt{cp}, \texttt{who}, \texttt{grep}, \texttt{gdb},
        and \texttt{mail} by system-program category.
  \item Why must \texttt{cd} be a shell built-in while \texttt{rm} can be a
        separate program?
  \item For each of the following, name the mechanism and the policy: (a) the
        timer that preempts a running process; (b) the permission check on
        \texttt{open}.
  \item In our counting model, a request passes through a chain of $k=3$
        servers and back. How many message traps and how many crossings does it
        cost, and how many times the monolithic cost is that?
  \item State the blast radius of the buggy driver under a modular kernel and
        under a virtual machine, and explain the difference.
\end{enumerate}

\subsection*{Solutions}

\begingroup\small
\begin{enumerate}
  \item \texttt{cp}: file management; \texttt{who}: status information;
        \texttt{grep}: file modification; \texttt{gdb}: programming-language
        support; \texttt{mail}: communications.
  \item A separate program runs in a child created by \texttt{fork};
        \texttt{chdir} there changes only the child's directory, and the child
        exits, so the shell's directory would not change. \texttt{rm} only
        needs to remove a file, which any process may request, so it can be a
        separate file that the shell loads.
  \item (a) Mechanism: the timer interrupt after a set interval; policy: how
        long the interval is. (b) Mechanism: the kernel's check of a permission
        bit before \texttt{open}; policy: who is granted which bits.
  \item $2k = 6$ message traps, $4k = 12$ crossings, versus 2 crossings for a
        monolithic call: $12/2 = 6$ times as many.
  \item Modular: the module runs in kernel mode, so the \texttt{NULL}
        dereference crashes the kernel --- the whole machine. Virtual machine:
        the driver runs in a guest kernel, so only that VM is lost; the manager
        and the other VMs keep running. Modules buy flexibility; virtual
        machines buy isolation.
\end{enumerate}
\endgroup

\begin{thebibliography}{9}
\bibitem{SilberschatzOSC_operating_system_concepts}
A.~Silberschatz, P.~B.~Galvin, and G.~Gagne, \emph{Operating System Concepts},
8th ed., John Wiley and Sons --- Ch.~2 \S2.5 (p.~66), \S2.6, \S2.7
(pp.~70--76), \S2.8 (p.~76). \emph{(Page-cited; index built 2026-10-14.)}
\bibitem{StallingsOS_internals_design_principles}
W.~Stallings, \emph{Operating Systems: Internals and Design Principles}, 6th
ed., Prentice Hall --- Ch.~2 \S2.3 (p.~64) and Ch.~4 \S4.3 (p.~179).
\emph{(Page-cited; index built 2026-10-14.)}
\bibitem{BovetCesati_understanding_linux_kernel}
D.~P.~Bovet and M.~Cesati, \emph{Understanding the Linux Kernel}, 3rd ed.,
O'Reilly Media, 2006 --- Ch.~1 (pp.~1--34). \emph{(Page-cited; index built
2026-10-14.)}
\end{thebibliography}

\end{document}
